\documentclass[11pt]{article}
\usepackage[margin=1in]{geometry}
\usepackage{amsmath,amssymb}
\usepackage{microtype}
\title{A Small Demonstration of Local Research Workflows}
\author{Octave Example Workspace}

\begin{document}
\maketitle

\begin{abstract}
This short paper exists to make Octave's editor, outline, citation audit, model context, and PDF workflow immediately testable after cloning the repository.
\end{abstract}

\section{Motivation}
Research software is most useful when the document remains the source of truth. Octave keeps analysis and conversation adjacent to ordinary files instead of importing them into a proprietary project format.

\section{A bounded context model}
Let $D$ denote the active document and let $P = \{p_1,\ldots,p_n\}$ be a finite set of explicitly pinned files. A model request receives
\[
  C = D \cup \bigcup_{i=1}^{n} p_i,
\]
subject to per-file and total context budgets. The model therefore sees deliberate project context, not an unbounded directory scrape.

\subsection{Review before writing}
Suggested revisions are represented as selectable diff hunks. The document is not modified until the researcher applies the selected proposal.

\section{Related principles}
The workflow follows the broader local-first principle that users should retain ownership and control of their data \cite{localfirst}.

\section{Conclusion}
The example is intentionally compact, but it exercises the same code paths used by a real research workspace.

\bibliographystyle{plain}
\bibliography{references}
\end{document}
